% ---------- IM244AI Operations Research (CIE - I, 02 April 2025)
\begin{iempaper}
\paperinfo{dept={INDUSTRIAL ENGINEERING \& MANAGEMENT},date={02\textsuperscript{nd} April 2025},exam={CIE -- I},maxmarks={10 + 50},sem={IV},prog={UG},duration={30 + 90 Min},
 title={OPERATIONS RESEARCH},code={IM244AI}}
\iemhead
\begin{cietable}{M}{BT}{CO}
\qpart{Part -- A}
\q{1}{The variables that represent decisions to be made in an optimization model are called \blank[1.8cm].}{1}{1}{1}
\q{2}{In an optimization model, the function that needs to be maximized or minimized is called the \blank[1.8cm].}{1}{1}{1}
\q{3}{Linear Programming is a mathematical technique used for \blank[1.8cm] decision-making.}{1}{2}{2}
\q{4}{One of the key limitations of Operations Research is that it requires \blank[1.8cm] data for accuracy.}{1}{2}{2}
\q{5}{The feasible region is the set of all possible solutions that satisfy all \blank[1.8cm].}{1}{1}{1}
\q{6}{The graphical solution method of LPP is applicable only when there are \blank[1.8cm] decision variables.}{1}{2}{2}
\q{7}{The additivity assumption ensures that the total effect of different decision variables is the \blank[1.8cm] of their individual effects.}{1}{2}{1}
\q{8}{When the total supply does not equal total demand, the transportation problem is called a(n) \blank[1.8cm] transportation problem.}{1}{2}{2}
\q{9}{A prohibited route transportation problem occurs when certain routes have \blank[1.8cm] restrictions and cannot be used for transportation.}{1}{2}{3}
\q{10}{The General Mathematical Formulation of a transportation problem consists of an objective function, supply constraints, and \blank[1.8cm] constraints.}{1}{2}{2}
\qpart{Part -- B}
\q{1}{Describe the Seven-Step Model-Building Process in Operations Research with an example.}{10}{2}{1}
\q{2}{Explain the key features of Operations Research (OR) models and highlight its limitations.}{10}{2}{1}
\q{3}{Consider the following problem faced by a production planner of a soft drink plant. He has two bottling machines A and B. A is designed for 8-ounce bottles and B for 16-ounce bottles. However, each can also be used for both types of bottles with some loss of efficiency. The manufacturing data is as follows:
\qtabc{|c|c|c|}{\hline Machine & 8-ounce Bottles & 16-ounce Bottles\\hline A & 100/minute & 40/minute\\hline B & 60/minute & 75/minute\\hline}
The machines can be run for 8 hours per day, 5 days per week. The profit on an 8-ounce bottle is Rs 1.5 and on a 16-ounce bottle is Rs 2.5. Weekly production of the drink cannot exceed 3,00,000 bottles and the market can absorb 25,000 8-ounce bottles and 7,000 16-ounce bottles per week. The planner wishes to maximize his profit, subject of course, to all the production and marketing restrictions. Formulate this problem as an LP model to maximize total profit.}{10}{3}{2}
\q{4}{A diet for a sick person must contain at least 4,000 units of vitamins, 50 units of minerals and 1,400 calories. Two foods A and B are available at a cost of Rs. 4 and Rs. 3 per unit, respectively. If one unit of A contains 200 units of vitamins, 1 unit of mineral and 40 calories and one unit of food B contains 100 units of vitamins, 2 units of minerals and 40 calories, formulate this problem as an LP model and solve it by graphical method to find the combination of foods to be used to have least cost.}{10}{3}{2}
\q{5}{Determine an initial basic feasible solution to the following transportation problem by using (a) NWCR, (b) LCM and (c) VAM.
\qtabc{|c|c|c|c|c|c|}{\hline & D1 & D2 & D3 & D4 & Supply\\hline S1 & 19 & 30 & 50 & 10 & 7\\hline S2 & 70 & 30 & 40 & 60 & 9\\hline S3 & 40 & 8 & 70 & 20 & 18\\hline Demand & 5 & 8 & 7 & 14 & \\hline}}{10}{3}{3}
\end{cietable}
\end{iempaper}
